% =========================================================
\chapter{Lecture 9 : Adversarial thinking II}
% =========================================================

\begin{rembox}[title=In a nutshell]
\textbf{What is this lecture about?} Lecture 8 asked ``what can go wrong?''. This one asks
\textbf{``what not to do?''}: the \textbf{Common Weakness Enumeration (CWE)}, MITRE's database of
the software errors that most often lead to vulnerabilities, grouped in three classes:
\textbf{insecure interaction between components}, \textbf{risky resource management} and
\textbf{porous defenses}. Most examples are on the web, so the lecture starts with a short
\textbf{web cheat sheet} (HTTP, cookies, GET and POST, HTML, PHP). We then dissect three classic
injection-style attacks: \textbf{OS command injection}, \textbf{cross-site scripting (XSS)} and
\textbf{cross-site request forgery (CSRF)}, with the \textbf{Same Origin Policy} and the defenses
against CSRF. \\
\textbf{Why is it useful?} XSS, SQL injection, CSRF and command injection sit in the top 10 of the
most dangerous software errors year after year. They all come from one mistake, using input you do
not control without sanitizing it, and one mental model, the confused deputy, explains CSRF.
\end{rembox}

\section{Web preliminaries}

\begin{rembox}[title=Why this cheat sheet]
Not a guide to web programming: only the basic concepts needed to understand the security problems
of the lecture. Up to now the course had two scenarios: a user acting \textbf{locally} on their
own device (local authentication, access control and execution), and a user acting
\textbf{remotely} on a server (the device is merely a screen). On the web, \textbf{browser and
server collaborate}: remote authentication, remote access control, and \textbf{mixed program
execution} (code runs on both sides). Two tools: \textbf{HTTP}, the protocol between browsers and
servers, and \textbf{PHP}, a language embedding dynamic, server-side elements in web pages.
\end{rembox}

\begin{defnbox}[title={Where we are: the top layers of the Internet (OSI model)}]
\begin{center}
\small
\begin{tabular}{l|l}
\textbf{Layer} & \textbf{Examples} \\ \hline
Application  & \textbf{HTTP} \\
Presentation & \\
Session      & SSL, TLS \\
Transport    & TCP, UDP \\
Network      & IP (naming and routing: DNS, BGP) \\
Data link    & IEEE 802.3 Ethernet (naming and routing: ARP) \\
Physical     & modulation and coding
\end{tabular}
\end{center}
This lecture focuses on the \textbf{top layers}.
\end{defnbox}

\begin{defnbox}[title={HTTP/1.1: HyperText Transfer Protocol}]
The protocol that determines what actions web servers and browsers take in response to commands.
\begin{itemize}
  \item \textbf{Request-response}: (1) the client (browser) sends a \textbf{request} (for an HTML
        file, to update a database, to send a mail, \ldots); (2) the server processes it, performs
        the action, and sends a \textbf{response} (the content to render, or an acknowledgement).
  \item \textbf{Stateless}: each command is executed independently, without any knowledge of the
        previous ones. Neither browser nor server keeps state: \textbf{the state must be sent
        inside the requests and responses}. Very important for security.
\end{itemize}
\end{defnbox}

\begin{defnbox}[title=Cookies]
A \textbf{small piece of data stored by the browser} on the user's device.
\begin{itemize}
  \item \textbf{Main goal}: store state information (shopping cart, login) to create HTTP
        \textbf{sessions} on top of a stateless protocol.
  \item \textbf{Secondary use}: tracking users.
  \item \textbf{Ambient authority in cookies}: if you are logged into \texttt{bank.com}, you have
        cookies for \texttt{bank.com}, and \textbf{every} new HTTP request to \texttt{bank.com}
        automatically includes all of them, so that you stay logged in (no login for every
        request). Whoever makes your browser send a request to \texttt{bank.com} uses your
        session.
\end{itemize}
\end{defnbox}

\begin{defnbox}[title={URL, GET and POST}]
\begin{itemize}
  \item \textbf{URL} (Uniform Resource Locator): a standard way to reference a resource (text,
        page, script, image). It contains the \textbf{protocol}, the \textbf{host} (a domain, or an
        IP and a port) and the \textbf{path} of the resource on that host, e.g.\
        \texttt{https://www.mywebsite.com/apparel/skirt.php}.
  \item \textbf{GET}: requests an existing resource. Its parameters are \textbf{encoded in the
        URL} as key/value pairs, the \textbf{query string}, after \texttt{?} and separated by
        \texttt{\&}: \texttt{\ldots/skirt.php?sku=123\&lang=en\&sect=silk} passes \texttt{sku=123},
        \texttt{lang=en}, \texttt{sect=silk}.
  \item \textbf{POST}: creates or updates a resource on the server. The data travels in the
        \textbf{request body} (form fields \texttt{name1=value1\&name2=value2}, JSON, XML).
\end{itemize}
A GET request \textbf{should not} change any data, a POST typically does; but \textbf{nothing
enforces it}. A developer could build \texttt{GET /user/delete?id=123}: if a search engine crawls
that link, it deletes a user. Very bad design, technically possible. (Standard methodologies, e.g.\
REST, structure HTTP APIs.)
\end{defnbox}

\begin{defnbox}[title={HTML and PHP}]
\begin{itemize}
  \item \textbf{HTML} (HyperText Markup Language): tells the browser \textbf{how to render} a
        document; parts of it are \textbf{marked with tags} (\texttt{<head>}, \texttt{<title>},
        \texttt{<body>}, \texttt{<h1>}, \texttt{<p>}, \ldots). The HTTP response usually carries
        HTML: data + the way to display it.
  \item \textbf{PHP} (``PHP Hypertext Preprocessor''): a \textbf{server} scripting language
        (still behind WordPress, Drupal) that builds web pages \textbf{on the fly} from inputs
        and variables. Variables start with \texttt{\$}. Special variables read what the client
        sent: \texttt{\$\_GET[param]} (from the URL), \texttt{\$\_POST[param]} (from the body),
        \texttt{\$\_SESSION[param]} (from the session tied to the cookie). \texttt{echo} outputs
        HTML.
\end{itemize}
\begin{lstlisting}[language=PHP]
<?php
$var = "class";
echo "<h2>PHP is Fun!</h2>";
echo "Hello $var!<br>";       // sent to the browser as: Hello class!<br>
echo "Learning PHP<br>";
?>
\end{lstlisting}
\end{defnbox}

\begin{defnbox}[title={Lifetime of a request on a PHP server}]
\textbf{GET}: the browser sends \texttt{GET /index.php} with headers (User-Agent, Accept, \ldots);
the web server (Apache, Nginx) sees the \texttt{.php} extension (not a static file) and passes the
request to the \textbf{PHP interpreter}; PHP runs the script command by command (reads the request,
the database, checks sessions, \ldots) and \textbf{builds an HTML document as a string}; the
complete HTML goes back to the web server and to the browser. \\
\textbf{POST}: same path, with the form data in the body (\texttt{username=test\&pass=1234}); PHP
fills \texttt{\$\_POST} and runs, e.g., \texttt{login.php}, typically writing to the database (the
main goal of a POST). It then answers with a page (\texttt{200 OK}, ``Success!'') or with a
\textbf{redirect} (\texttt{header('Location: /dashboard.php')}): the
\textbf{Post/Redirect/Get} pattern, which avoids submitting the form twice when the user hits
``refresh''.
\end{defnbox}

\section{Common Weakness Enumeration (CWE)}

\begin{defnbox}[title={Common Weakness Enumeration}]
\textbf{Idea}: a \textbf{database of software errors leading to vulnerabilities}, to help security
engineers avoid common pitfalls: \textbf{``what not to do''}. STRIDE (Lecture 8) reasons about what
the adversary can do; the CWE reduces the attack surface by \textbf{not repeating known errors}.
MITRE explains each error with its consequences; at a high level, each one makes the software
break a security principle. The 2011 classification has three classes:
\begin{itemize}
  \item \textbf{Insecure interaction between components}: one subsystem feeds another subsystem
        data that is not sanitized.
  \item \textbf{Risky resource management}: the system acts on inputs that are not sanitized.
  \item \textbf{Porous defenses}: defenses fail to provide full protection or complete mediation,
        through missing checks or partial mechanisms.
\end{itemize}
\end{defnbox}

\begin{exbox}[title={CWE Top 25 (current list, first 8)}]
\begin{center}
\small
\begin{tabular}{r|l|l}
\textbf{\#} & \textbf{Weakness} & \textbf{CWE} \\ \hline
1 & Cross-site Scripting (XSS) & CWE-79 \\
2 & Out-of-bounds Write & CWE-787 \\
3 & SQL Injection & CWE-89 \\
4 & Cross-Site Request Forgery (CSRF) & CWE-352 \\
5 & Path Traversal & CWE-22 \\
6 & Out-of-bounds Read & CWE-125 \\
7 & OS Command Injection & CWE-78 \\
8 & Use After Free & CWE-416
\end{tabular}
\end{center}
XSS, SQL injection, CSRF and OS command injection (all \emph{injections} of attacker-controlled
input) are highlighted in the slides: four of the top 7.
\end{exbox}

\section{CWE I: Insecure interaction between components}

\begin{defnbox}[title={Insecure interaction between components}]
``Insecure ways in which data is sent and received between separate components, modules, programs,
processes, threads, or systems.'' One subsystem feeds another subsystem data that is \textbf{not
sanitized}; the program uses it, with unintended behaviors that the adversary exploits.
\end{defnbox}

\subsection{CWE-78: OS command injection}

\begin{defnbox}[title={CWE-78: Improper Neutralization of Special Elements used in an OS Command}]
A program shows the user the content of their folder under \texttt{/home}. A web form with one field
sends \texttt{userName} by POST; the server pastes it at the end of the Linux command
\texttt{ls -l /home/} (for \texttt{ctroncoso}: \texttt{ls -l /home/ctroncoso}).
\begin{lstlisting}[language=HTML]
<form action="/url/myscript.php" method="post">
  userName: <input name="userName" type="text" />
  <input name="submit" type="submit" value="Submit">
</form>
\end{lstlisting}
\begin{lstlisting}[language=PHP]
$userName = $_POST["userName"];       // no check on $userName format!
$command  = 'ls -l /home/' . $userName;
system($command);
\end{lstlisting}
\end{defnbox}
